\documentclass[11pt, a4paper]{article}
% ══════════════════════════════════════════════════════════════
% Preâmbulo compartilhado pelos documentos da P2
% Reaproveita o estilo de docs/documento-unificado.tex
% ══════════════════════════════════════════════════════════════

\usepackage[utf8]{inputenc}
\usepackage[T1]{fontenc}
\usepackage[brazil]{babel}

\usepackage[top=2.5cm, bottom=2.5cm, left=3cm, right=2cm]{geometry}
\usepackage{setspace}
\onehalfspacing
\setlength{\parindent}{0pt}
\setlength{\parskip}{6pt}

\usepackage{lmodern}
\usepackage{microtype}
\usepackage{amssymb}
\usepackage{amsmath}
\usepackage{pifont}
\newcommand{\cmark}{\ding{51}}
\newcommand{\xmark}{\ding{55}}

\usepackage[table,dvipsnames]{xcolor}
\definecolor{accent}{HTML}{2C7A4B}
\definecolor{accentdim}{HTML}{1E5934}
\definecolor{accentlight}{HTML}{E8F4ED}
\definecolor{tabheader}{HTML}{2C3E50}
\definecolor{tabrow}{HTML}{F5F6FA}
\definecolor{linkblue}{HTML}{2980B9}
\definecolor{codebg}{HTML}{F5F5F5}
\definecolor{codeborder}{HTML}{D0D0D0}
\definecolor{ndvilow}{HTML}{8B4513}
\definecolor{ndvimid}{HTML}{DAA520}
\definecolor{ndvihigh}{HTML}{006400}
\definecolor{warnyellow}{HTML}{F39C12}
\definecolor{successgreen}{HTML}{27AE60}
\definecolor{errorred}{HTML}{E74C3C}

\usepackage{booktabs}
\usepackage{tabularx}
\usepackage{multirow}
\usepackage{makecell}
\usepackage{float}
\usepackage{graphicx}
\usepackage{caption}
\usepackage{array}

\usepackage{enumitem}
\setlist[itemize]{itemsep=2pt, parsep=0pt, topsep=4pt, leftmargin=*}
\setlist[enumerate]{itemsep=2pt, parsep=0pt, topsep=4pt, leftmargin=*}

\usepackage[colorlinks=true, linkcolor=linkblue, urlcolor=linkblue, citecolor=linkblue]{hyperref}

\usepackage{fancyhdr}
\pagestyle{fancy}
\fancyhf{}
\fancyhead[L]{\small\textcolor{gray}{EasyGis}}
\fancyhead[R]{\small\textcolor{gray}{P2 — Gerência e Manutenção de Software · Prof.\ Mauro · 2026.1}}
\fancyfoot[C]{\thepage}
\renewcommand{\headrulewidth}{0.4pt}

\usepackage{titlesec}
\titleformat{\section}
  {\Large\bfseries\color{accentdim}}
  {\thesection.}{0.5em}{}
\titleformat{\subsection}
  {\large\bfseries\color{accent}}
  {\thesubsection}{0.5em}{}
\titleformat{\subsubsection}
  {\normalsize\bfseries}
  {\thesubsubsection}{0.5em}{}

% Listings (código)
\usepackage{listings}
\lstdefinestyle{shellstyle}{
  basicstyle=\ttfamily\small,
  backgroundcolor=\color{codebg},
  frame=single,
  rulecolor=\color{codeborder},
  framesep=4pt,
  xleftmargin=8pt,
  xrightmargin=4pt,
  breaklines=true,
  showstringspaces=false,
  keepspaces=true,
  literate={─}{{-}}1 {├}{{|}}1 {│}{{|}}1 {└}{{`}}1
}
\lstdefinestyle{pythonstyle}{
  basicstyle=\ttfamily\footnotesize,
  language=Python,
  backgroundcolor=\color{codebg},
  frame=single,
  rulecolor=\color{codeborder},
  framesep=4pt,
  xleftmargin=8pt,
  xrightmargin=4pt,
  breaklines=true,
  keywordstyle=\color{accent}\bfseries,
  stringstyle=\color{ndvimid},
  commentstyle=\color{gray}\itshape,
  showstringspaces=false,
  keepspaces=true
}
\lstset{style=shellstyle}

% Caixas para código inline
\newcommand{\code}[1]{\texttt{\colorbox{codebg}{\strut#1}}}

% Caixas de destaque
\usepackage{tcolorbox}
\tcbuselibrary{skins, breakable}

\newtcolorbox{infobox}[1][]{%
  colback=accentlight, colframe=accent,
  boxrule=0.6pt, arc=2pt,
  left=8pt, right=8pt, top=4pt, bottom=4pt,
  fonttitle=\bfseries, breakable, #1
}

\newtcolorbox{warnbox}[1][]{%
  colback=yellow!8, colframe=warnyellow,
  boxrule=0.6pt, arc=2pt,
  left=8pt, right=8pt, top=4pt, bottom=4pt,
  breakable, #1
}

\newtcolorbox{successbox}[1][]{%
  colback=green!4, colframe=successgreen,
  boxrule=0.6pt, arc=2pt,
  left=8pt, right=8pt, top=4pt, bottom=4pt,
  breakable, #1
}


\title{%
  \vspace{-1.5cm}
  {\Huge\bfseries\color{accentdim} Entrega da Avaliação P2}\\[8pt]
  {\Large\bfseries Engenharia de Software --- 2026.1}\\[18pt]
  {\large ICEV Remote Sensing Software}\\[4pt]
  {\normalsize\itshape Plataforma de Sensoriamento Remoto para Agricultura de Precisão}
}
\author{%
  Luis Gustavo Olimpio \and Lauan Matheus \and João Leonardi\\
  \and João Vinícius Castello \and Vinícius Henrique\\
  \and Lucas Benevinuto \and José Melquíades
}
\date{Maio de 2026}

\begin{document}
\maketitle
\thispagestyle{empty}

\vfill
\begin{center}
\small\textit{Instituto de Ensino Superior --- ICEV}\\
\small\textit{Teresina, PI, Brasil}\\[6pt]
\small\textit{Disciplina: Engenharia de Software --- Avaliação P2}\\
\small\textit{Prazo de entrega: 11/05/2026}
\end{center}
\newpage

\tableofcontents
\newpage

% ══════════════════════════════════════════════════════════════
\section{Mapeamento dos Entregáveis Obrigatórios}
% ══════════════════════════════════════════════════════════════

Este documento é o ponto de partida da avaliação. Cada entregável obrigatório da P2 está mapeado para o arquivo correspondente no repositório do projeto.

\renewcommand{\arraystretch}{1.4}
\begin{table}[H]
\centering
\small
\rowcolors{2}{tabrow}{white}
\begin{tabularx}{\textwidth}{@{}clX@{}}
\toprule
\rowcolor{tabheader}\textcolor{white}{\textbf{\#}} & \textcolor{white}{\textbf{Entregável da P2}} & \textcolor{white}{\textbf{Localização}}\\
\midrule
1 & Código-fonte do MVP & Repositório completo --- diretórios \texttt{api/}, \texttt{web/}, \texttt{data/} com versionamento Git\\
2 & Plano de Testes & \texttt{docs/p2/08-plano-de-testes.pdf}\\
2 & Relatório de Testes & \texttt{docs/p2/09-relatorio-de-testes.pdf}\\
3 & Manual do Usuário & \texttt{docs/p2/07-manual-usuario.pdf}\\
4 & Apresentação e Demonstração & \texttt{docs/p2/apresentacao-p2.pdf}\\
5 & Documentação Técnica & Ver Seção~\ref{sec:doctec} (consolidada)\\
\bottomrule
\end{tabularx}
\end{table}

% ══════════════════════════════════════════════════════════════
\section{Documentação Técnica Consolidada}
\label{sec:doctec}
% ══════════════════════════════════════════════════════════════

A documentação técnica está \textbf{distribuída em vários documentos especializados} já produzidos no Sprint inicial do projeto (TP1). O escopo de cada um é apresentado na Tabela~\ref{tab:doctec}.

\begin{table}[H]
\centering
\small
\rowcolors{2}{tabrow}{white}
\begin{tabularx}{\textwidth}{@{}lX@{}}
\toprule
\rowcolor{tabheader}\textcolor{white}{\textbf{Item exigido pela P2}} & \textcolor{white}{\textbf{Documento(s)}}\\
\midrule
Descrição detalhada da arquitetura & \texttt{docs/04-diagramas-arquitetura.md} --- 8 diagramas (alto nível, componentes, deployment, sequência NDVI, sequência classificação, estrutura de diretórios, fluxo de dados)\\
Tecnologias utilizadas & \texttt{docs/04-diagramas-arquitetura.md} \S 8 --- stack consolidada frontend + backend\\
Decisões de design e justificativas & \texttt{docs/04-diagramas-arquitetura.md} \S 9 (DA-01 a DA-05) + complemento na Seção~\ref{sec:decisoes} deste documento (DA-06 a DA-11)\\
Requisitos funcionais e não funcionais & \texttt{docs/03-especificacao-requisitos-SRS.md} --- 20 RFs e 8 RNFs\\
Backlog priorizado e rastreabilidade & \texttt{docs/02-backlog-produto.md} --- 33 histórias em 6 épicos\\
Visão de produto e usuários-alvo & \texttt{docs/01-documento-visao-produto.md}\\
Protótipo de interface & \texttt{docs/05-prototipo-interface.md}\\
Plano de projeto & \texttt{docs/06-plano-projeto-inicial.md}\\
\bottomrule
\end{tabularx}
\caption{Mapeamento dos itens da documentação técnica}
\label{tab:doctec}
\end{table}

\begin{infobox}
Os documentos do TP1 também estão consolidados em \texttt{docs/tp1/documento-unificado.pdf}, gerado a partir de \texttt{docs/tp1/documento-unificado.tex}.
\end{infobox}

% ══════════════════════════════════════════════════════════════
\section{Decisões de Design Complementares}
\label{sec:decisoes}
% ══════════════════════════════════════════════════════════════

Em complemento às decisões DA-01 a DA-05 já documentadas em \texttt{docs/04-diagramas-arquitetura.md} \S 9, estas são as decisões adicionais relevantes para a entrega do MVP.

\subsection*{DA-06 --- Separação \texttt{SentinelProcessor} \(\leftrightarrow\) Camada de Rotas}

\textbf{Decisão.} A classe \texttt{SentinelProcessor} (em \texttt{api/sentinel\_processor.py}) concentra toda a lógica de I/O geoespacial e cálculo numérico. As rotas FastAPI em \texttt{api/main.py} apenas orquestram e serializam.

\textbf{Justificativa.} Permite testar a camada de processamento em isolamento sem precisar subir o servidor HTTP --- exatamente o que fazemos em \texttt{tests/test\_sentinel\_processor.py}. Reduz acoplamento e melhora a manutenibilidade (RNF-05).

\subsection*{DA-07 --- Tipagem Estática End-to-End}

\textbf{Decisão.} TypeScript estrito no frontend e Pydantic no backend, com modelos espelhados (\texttt{Coordinate}, \texttt{IndexResult}, \texttt{ClassificationResult}).

\textbf{Justificativa.} Erros de contrato entre cliente e servidor são detectados em build/runtime, não em produção. O Pydantic também valida automaticamente a entrada do usuário, eliminando código manual de validação.

\subsection*{DA-08 --- Resolução do Produto Sentinel-2 mais Recente como Padrão}

\textbf{Decisão.} Quando \texttt{product\_name} é omitido, o backend ordena os produtos \texttt{.SAFE} por nome (que contém data ISO) e usa o mais recente.

\textbf{Justificativa.} Otimiza o caminho feliz do usuário leigo, que normalmente quer ``ver agora''. O usuário avançado ainda pode passar um produto específico para comparações temporais.

\subsection*{DA-09 --- Suavização Gaussiana com Preservação de NaN}

\textbf{Decisão.} O filtro gaussiano (\texttt{apply\_smooth\_filter}) divide o resultado pela máscara também filtrada, preservando o limite do polígono do talhão.

\textbf{Justificativa.} Aplicar gaussiano direto ``sangra'' os valores válidos para fora do talhão, criando halos artificiais. A correção por máscara é a abordagem padrão em sensoriamento remoto.

\subsection*{DA-10 --- Visualização 3D via Downsampling}

\textbf{Decisão.} Para o terreno 3D (\texttt{web/components/terrain-3d-viewer.tsx}), reduz-se a resolução por fator 4 antes de enviar ao cliente.

\textbf{Justificativa.} Talhões grandes podem ter milhões de pixels. Renderizar isso direto em WebGL trava navegadores em GPU integrada (RNF-01). O downsample mantém a interpretação visual sem custo de performance.

\subsection*{DA-11 --- Gerenciamento de Dependências Python via \texttt{uv}}

\textbf{Decisão.} Usar \texttt{uv} (\texttt{pyproject.toml + uv.lock}) em vez de \texttt{pip + requirements.txt}.

\textbf{Justificativa.} O \texttt{uv} é cerca de dez vezes mais rápido em instalação, garante builds reprodutíveis via \textit{lock file} e simplifica o setup em máquinas novas (\texttt{uv sync} resolve tudo). \textit{Trade-off:} dependência extra que o usuário precisa instalar --- mitigado pela documentação clara no \texttt{README.md}.

% ══════════════════════════════════════════════════════════════
\section{Tecnologias --- Justificativa Resumida}
% ══════════════════════════════════════════════════════════════

A tabela completa está em \texttt{docs/04-diagramas-arquitetura.md} \S 8. As justificativas-chave estão na Tabela~\ref{tab:tech}.

\begin{table}[H]
\centering
\small
\rowcolors{2}{tabrow}{white}
\begin{tabularx}{\textwidth}{@{}lX@{}}
\toprule
\rowcolor{tabheader}\textcolor{white}{\textbf{Tecnologia}} & \textcolor{white}{\textbf{Por que esta e não outra}}\\
\midrule
FastAPI & Validação automática via Pydantic, OpenAPI grátis (\texttt{/docs}), \textit{async} nativo. Alternativas (Flask) exigiriam plugins para o mesmo nível.\\
rasterio & Padrão de fato para I/O de imagens geoespaciais em Python. Wrapper Pythonic sobre GDAL.\\
Next.js + React & App Router permite SSR + API routes no mesmo projeto (ex.: \texttt{/api/fields} lê KML do disco sem expor o filesystem).\\
Leaflet & Maduro, leve, integra bem com tiles do Esri/OSM. Alternativas (Mapbox GL) exigem token e são pagas.\\
Three.js / react-three-fiber & Visualização 3D nativa em WebGL. Alternativa (Cesium) é \textit{overkill} para terreno simples.\\
Tailwind + shadcn/ui & Componentes acessíveis (Radix por baixo), sem CSS-in-JS, sem bundler extra.\\
fast-xml-parser & KML é XML; este parser é o mais rápido em Node.js sem dependências nativas.\\
\bottomrule
\end{tabularx}
\caption{Justificativas de escolha de tecnologias}
\label{tab:tech}
\end{table}

% ══════════════════════════════════════════════════════════════
\section{Como Avaliar Esta Entrega}
% ══════════════════════════════════════════════════════════════

\subsection{Setup rápido (5 minutos)}

\begin{lstlisting}[style=shellstyle]
# Pre-requisitos: Python 3.12+, Node 18+, GDAL (brew install gdal)
uv sync                       # instala dependencias Python
cd web && npm install         # instala dependencias Node
cd ..

# Em 2 terminais:
./run_api.sh                  # backend na porta 8000
cd web && npm run dev         # frontend na porta 3000
\end{lstlisting}

\subsection{Roteiro de demonstração (5 minutos)}

\begin{enumerate}
  \item Abrir \url{http://localhost:3000}
  \item Selecionar talhão \texttt{crop field 1}
  \item Calcular NDVI \(\rightarrow\) ver overlay no mapa, estatísticas e histograma
  \item Trocar para EVI \(\rightarrow\) comparar
  \item Ativar visualização 3D \(\rightarrow\) rotacionar terreno
  \item Ir em \texttt{/classification} \(\rightarrow\) rodar K-Means com 3 classes \(\rightarrow\) ver áreas em hectares
\end{enumerate}

\subsection{Para executar a suíte de testes}

\begin{lstlisting}[style=shellstyle]
uv sync --group dev           # instala pytest e dependencias de teste
uv run pytest tests/ -v       # roda 28 testes - ver relatorio em docs/09
\end{lstlisting}

% ══════════════════════════════════════════════════════════════
\section{Estrutura Final do Repositório}
% ══════════════════════════════════════════════════════════════

\begin{lstlisting}[style=shellstyle]
icev-remote-sensing-software/
|-- README.md                  # Setup + instrucoes de execucao
|-- pyproject.toml             # Dependencias Python
|-- api/                       # Backend FastAPI
|   |-- main.py
|   `-- sentinel_processor.py
|-- web/                       # Frontend Next.js
|   |-- app/
|   |-- components/
|   |-- lib/
|   `-- types/
|-- data/
|   |-- products/              # Imagens Sentinel-2 .SAFE
|   `-- KML Fields/            # Poligonos de talhoes
|-- tests/                     # Suite de testes (pytest)  [NOVO P2]
|   |-- conftest.py
|   |-- test_sentinel_processor.py
|   `-- test_api_endpoints.py
`-- docs/
    |-- README.md                  # Indice geral
    |-- 01-documento-visao-produto.md      # TP1 - referenciados pela P2
    |-- 02-backlog-produto.md
    |-- 03-especificacao-requisitos-SRS.md
    |-- 04-diagramas-arquitetura.md
    |-- 05-prototipo-interface.md
    |-- 06-plano-projeto-inicial.md
    |-- p2/                        # [NOVO] Entrega P2
    |   |-- 00-ENTREGA-P2.{tex,pdf}     # <- Voce esta aqui
    |   |-- 07-manual-usuario.{tex,pdf}
    |   |-- 08-plano-de-testes.{tex,pdf}
    |   |-- 09-relatorio-de-testes.{tex,pdf}
    |   |-- apresentacao-p2.{tex,pdf,pptx}
    |   |-- p2-preamble.tex
    |   |-- Makefile
    |   |-- build_slides.py
    |   `-- md-sources/            # Originais markdown
    `-- tp1/                       # Material complementar do TP1
        |-- documento-unificado.{tex,pdf}
        |-- canvas-experimento.tex
        |-- canvas-experimentofinal.pdf
        |-- atividade-unidade4.{tex,pdf}
        `-- images/
\end{lstlisting}

\end{document}
